%% ma: L12-B14-0048
%% lop: 12
%% chuong: 5
%% bai: 14
%% dang: 12.14.1
%% loai: TN4
%% muc: NB
%% dap_an: "A"
%% nguon: "To_n_12_PT_M_t_Ph_ng.pdf (Đề kiểm tra đánh giá kiến thức, Chuyên đề Phương trình mặt phẳng, Nguyễn Anh Khôi - Lê Minh Khôi), Câu 2"
%% kien_thuc: "Bài 14 (phương trình tổng quát của mặt phẳng, vectơ pháp tuyến)"
%% trang_thai: chinh-thuc
%% ghi_chu: "Tài liệu gốc không kèm đáp án; đáp án do Claude giải và kiểm bằng sympy."
\begin{ex}
Trong không gian $Oxyz$, mặt phẳng $(P)\colon 2x-z+3=0$ có một vectơ pháp tuyến là
\choice{\True $\vec n_1=(2;0;-1)$}{$\vec n_2=(2;-1;3)$}{$\vec n_3=(2;-1;0)$}{$\vec n_4=(-1;0;-1)$}
\loigiai{$(P)\colon 2x+0y-z+3=0$ có vectơ pháp tuyến $\vec n_1=(2;0;-1)$. Chọn A.}
\end{ex}
